% BEGIN REV09_FRICKE_JET
\subsection{Correspondencia de Fricke y recuperación diferencial}
\label{corpus9:iv:fricke}

La traza reticular de nivel dos y su completación de Fricke
proporcionan, después de la construcción excepcional, las
coordenadas
\[
 t=t_2,\qquad s=\frac{4096}{t},\qquad
 T=t+s+24,\qquad t\in\mathbb C^\times.
\]
La involución intercambia \(t\) y \(s\). Las dos lecturas modulares
ordenadas son
\[
 j_1=\frac{(t+256)^3}{t^2},\qquad
 j_2=\frac{(s+256)^3}{s^2}.
\]
Estas fórmulas utilizan la uniformización ya establecida en
\eqref{km:j-nivel2} y conservan la elección de la lectura orientada.

\begin{proposition}[Correspondencia y discriminante]
\label{corpus9:iv:fricke-correspondencia}
Se cumplen las identidades
\begin{align}
 j_1-j_2&=(s-t)(T+23),\\
 j_1+j_2&=T^2+T-7256,\\
 j_1j_2&=(T+248)^3.
 \label{corpus9:iv:fricke-simetricos}
\end{align}
Los dos valores son raíces de
\[
 X^2-(T^2+T-7256)X+(T+248)^3=0,
\]
cuyo discriminante es
\begin{equation}
 \Delta=(T-152)(T+104)(T+23)^2.
 \label{corpus9:iv:fricke-discriminante}
\end{equation}
\end{proposition}
\begin{proof}
La expansión de la primera lectura, usando \(ts=4096\), da
\(j_1=t+768+48s+s^2\); la segunda se obtiene intercambiando
\(s,t\). La resta da \((s-t)(t+s+47)\), y la suma se reduce a
\((t+s)^2+49(t+s)-6656\). Sustituir \(t+s=T-24\) produce
las dos primeras identidades. Para el producto,
\[
 j_1j_2
 =\frac{((t+256)(s+256))^3}{(ts)^2}
 =(T+248)^3.
\]
Finalmente,
\((s-t)^2=(T-24)^2-16384=(T-152)(T+104)\).
Multiplicar por \((T+23)^2\) prueba
\eqref{corpus9:iv:fricke-discriminante}.
\end{proof}

\begin{proposition}[Recuperación de la orientación por valores y primer jet]
\label{corpus9:iv:fricke-jet}
Para \(T\ne-23\), la pareja ordenada \((T,j_1)\) recupera
\begin{equation}
 t=\frac{T^2-j_1-3904}{T+23}.
 \label{corpus9:iv:fricke-inversa}
\end{equation}
En \(T=-23\), los dos parámetros
\[
 t_\pm=\frac{-47\pm45i\sqrt7}{2}
\]
comparten \((T,j_1)=(-23,-3375)\). El primer jet
\(p=dj_1/dT\), calculado en la rama parametrizada por \(t\),
los separa y permite recuperar
\begin{equation}
 p_\pm=\frac{-45\mp45i\sqrt7}{2},\qquad
 s=p-1,\qquad t=-46-p.
 \label{corpus9:iv:fricke-jet-inversa}
\end{equation}
\end{proposition}
\begin{proof}
De \(s^2=(T-24)s-4096\) resulta
\[
 j_1=(T+23)s+T-3352
     =T^2-3904-(T+23)t.
\]
Para \(T\ne-23\), esta igualdad da
\eqref{corpus9:iv:fricke-inversa}. Para \(T=-23\),
las condiciones \(t+s=-47\), \(ts=4096\) dan
\(t^2+47t+4096=0\), de discriminante \(-45^2\cdot7\).
La misma identidad da \(j_1=-3375\) en ambas raíces.

La derivada \(dT/dt=1-4096/t^2\) sólo se anula en
\(t=\pm64\). Ninguna de las dos raíces excepcionales tiene ese
valor, por lo que \(T\) es coordenada local en ambas ramas.
Diferenciando la primera expresión de \(j_1\),
\[
 \frac{dj_1}{dt}
  =(s+1)\frac{dT}{dt}+(T+23)\frac{ds}{dt}.
\]
En \(T=-23\), resulta \(p=s+1=-46-t\), que prueba
\eqref{corpus9:iv:fricke-jet-inversa}, con signos opuestos
para los índices emparejados \(t_\pm,p_\pm\).
\end{proof}

Los puntos fijos de Fricke, \(t=\pm64\), expresan respectivamente
\((T,j)=(152,8000)\) y \((-104,1728)\). La coincidencia en
\(T=-23\) corresponde en cambio a dos puntos distintos de la
parametrización. El campo de definición se obtiene de la ecuación:
\[
 \mathbb Q(t_\pm)=\mathbb Q(\sqrt{-7}),\qquad
 \left(\frac{2t+47}{45}\right)^2=-7.
\]
Sobre esta fibra, Fricke coincide con la conjugación cuadrática,
de traza \(-47\) y norma \(4096\). El jet conserva ese campo,
pues \((2p+45)^2=-45^2\cdot7\). La reconstrucción es exacta
en esta correspondencia modular: el dato diferencial distingue
las dos orientaciones que comparten sus valores escalares.
% END REV09_FRICKE_JET
